% GENERATED FILE. Edit canonical lessons, then run npm run generate:books.
% canonical-source-hash: fnv1a64:ad1b6e0f1248aed6
% canonical-lessons: RU-C325-imperiya, RU-C325-revolyutsiya, RU-C325-vosstaniye, RU-C325-zavoyovyvat, RU-C325-pravit

\chapter{Empires and uprisings}
\label{ch:ru-empires-and-uprisings}
\hlchaptermodality{325}

\begin{chapteropening}
\textbf{By the end of this chapter:} \emph{I can talk about an empire, a revolution and an uprising, and say that someone conquered or ruled a land.}

Complete the last lesson of chapter 325: I can talk about an empire, a revolution and an uprising, and say that someone conquered or ruled a land.
\end{chapteropening}

\section[\texorpdfstring{\ru{империя} (impériya)}{империя (impériya)}]{\ru{империя} (impériya) — an empire}

\label{lesson:RU-C325-imperiya}



\begin{quote}
Before the new one: say the Russian for to die, then the Russian for to free.
\end{quote}

\subsection*{You'll want to know: \ru{империя}}
\textbf{\ru{империя}} (\emph{impériya}) — \textquotedblleft{}an empire\textquotedblright{}.

\begin{quote}
\ru{Империя} \ru{существовала} \ru{тысячу} \ru{лет}.
\end{quote}

\begin{quote}
\emph{Impériya sushchestvovála týsyachu let.}
\end{quote}

\begin{quote}
The empire lasted a thousand years.
\end{quote}

Feminine.

\begin{cousinweb}
\emph{\ru{империя}} goes back to Latin \emph{imperium} \textquotedblleft{}command, rule\textquotedblright{} (EtymDB), from \emph{imperāre} \textquotedblleft{}to command\textquotedblright{}.
\end{cousinweb}

\begin{grammarlens}[title={What you've built}]
One more word for telling, explaining and describing at B1.
\end{grammarlens}

\subsection*{Your turn}
\begin{itemize}
\raggedright
  \item \emph{Say it:} \emph{impériya}
  \item \emph{Say it:} \emph{impériya}, once more
  \item \emph{Say it:} a sentence of your own with \emph{impériya}
  \item \emph{Recall:} say the Russian for to die, then the Russian for to free, then say \emph{impériya} again
\end{itemize}

\begin{tcolorbox}[breakable,colback=teal!4,colframe=teal!35!black,title={Before you move on}]
What does \ru{империя} mean? (\textquotedblleft{}An empire\textquotedblright{}.) Say it once more.
\end{tcolorbox}
\section[\texorpdfstring{\ru{революция} (revolyútsiya)}{революция (revolyútsiya)}]{\ru{революция} (revolyútsiya) — a revolution}

\label{lesson:RU-C325-revolyutsiya}



\begin{quote}
Before the new one: say the Russian for to free, then the Russian for an empire.
\end{quote}

\subsection*{You'll want to know: \ru{революция}}
\textbf{\ru{революция}} (\emph{revolyútsiya}) — \textquotedblleft{}a revolution\textquotedblright{}.

\begin{quote}
\ru{После} \ru{революции} \ru{жизнь} \ru{в} \ru{стране} \ru{изменилась}.
\end{quote}

\begin{quote}
\emph{Pósle revolyútsii zhizn v strané izmenílas.}
\end{quote}

\begin{quote}
After the revolution life in the country changed.
\end{quote}

Feminine.

\begin{cousinweb}
\emph{\ru{революция}} came through Polish \emph{rewolucja} from Latin \emph{revolūtiō} (EtymDB), \textquotedblleft{}a revolving\textquotedblright{} (Lewis \& Short), built on \emph{revolvere} \textquotedblleft{}to roll back\textquotedblright{}, from \emph{volvere} \textquotedblleft{}to roll\textquotedblright{}.
\end{cousinweb}

\begin{grammarlens}[title={What you've built}]
One more word for telling, explaining and describing at B1.
\end{grammarlens}

\subsection*{Your turn}
\begin{itemize}
\raggedright
  \item \emph{Say it:} \emph{revolyútsiya}
  \item \emph{Say it:} \emph{revolyútsiya}, once more
  \item \emph{Say it:} a sentence of your own with \emph{revolyútsiya}
  \item \emph{Recall:} say the Russian for to free, then the Russian for an empire, then say \emph{revolyútsiya} again
\end{itemize}

\begin{tcolorbox}[breakable,colback=teal!4,colframe=teal!35!black,title={Before you move on}]
What does \ru{революция} mean? (\textquotedblleft{}A revolution\textquotedblright{}.) Say it once more.
\end{tcolorbox}
\section[\texorpdfstring{\ru{восстание} (vosstániye)}{восстание (vosstániye)}]{\ru{восстание} (vosstániye) — an uprising, a revolt}

\label{lesson:RU-C325-vosstaniye}



\begin{quote}
Before the new one: say the Russian for an empire, then the Russian for a revolution.
\end{quote}

\subsection*{You'll want to know: \ru{восстание}}
\textbf{\ru{восстание}} (\emph{vosstániye}) — \textquotedblleft{}an uprising, a revolt\textquotedblright{}.

\begin{quote}
\ru{Восстание} \ru{началось} \ru{на} \ru{юге} \ru{страны}.
\end{quote}

\begin{quote}
\emph{Vosstániye nachalós na yúge straný.}
\end{quote}

\begin{quote}
The uprising began in the south of the country.
\end{quote}

Neuter.

\begin{cousinweb}
\emph{\ru{восстание}} is built on \emph{\ru{восстать}} \textquotedblleft{}to rise up\textquotedblright{}, from \emph{\ru{вос}-} \textquotedblleft{}up\textquotedblright{} and \emph{\ru{стать}} \textquotedblleft{}to stand\textquotedblright{}, inherited from Proto-Slavic \emph{stati}, PIE \emph{steh\textsubscript{2}-} (EtymDB).
\end{cousinweb}

\begin{grammarlens}[title={What you've built}]
One more word for telling, explaining and describing at B1.
\end{grammarlens}

\subsection*{Your turn}
\begin{itemize}
\raggedright
  \item \emph{Say it:} \emph{vosstániye}
  \item \emph{Say it:} \emph{vosstániye}, once more
  \item \emph{Say it:} a sentence of your own with \emph{vosstániye}
  \item \emph{Recall:} say the Russian for an empire, then the Russian for a revolution, then say \emph{vosstániye} again
\end{itemize}

\begin{tcolorbox}[breakable,colback=teal!4,colframe=teal!35!black,title={Before you move on}]
What does \ru{восстание} mean? (\textquotedblleft{}An uprising, a revolt\textquotedblright{}.) Say it once more.
\end{tcolorbox}
\section[\texorpdfstring{\ru{завоёвывать} (zavoyóvyvat)}{завоёвывать (zavoyóvyvat)}]{\ru{завоёвывать} (zavoyóvyvat) — to conquer}

\label{lesson:RU-C325-zavoyovyvat}



\begin{quote}
Before the new one: say the Russian for a revolution, then the Russian for an uprising.
\end{quote}

\subsection*{You'll want to know: \ru{завоёвывать}}
\textbf{\ru{завоёвывать}} (\emph{zavoyóvyvat}) — \textquotedblleft{}to conquer\textquotedblright{}.

\begin{quote}
\ru{Они} \ru{завоевали} \ru{много} \ru{земель}.
\end{quote}

\begin{quote}
\emph{Oní zavoyeváli mnógo zemél.}
\end{quote}

\begin{quote}
They conquered many lands.
\end{quote}

Imperfective; the perfective is \emph{\ru{завоевать}} (\emph{zavoyevát}), \emph{\ru{завоюю}}. Also \emph{\ru{завоевать} \ru{сердце}} — to win a heart.

\begin{cousinweb}
\emph{\ru{завоёвывать}} is built on \emph{\ru{за}-} and \emph{\ru{воевать}} \textquotedblleft{}to wage war\textquotedblright{}, from \emph{\ru{воин}} \textquotedblleft{}a warrior\textquotedblright{}, Proto-Slavic \emph{vojinŭ} (EtymDB).
\end{cousinweb}

\begin{grammarlens}[title={What you've built}]
One more word for telling, explaining and describing at B1.
\end{grammarlens}

\subsection*{Your turn}
\begin{itemize}
\raggedright
  \item \emph{Say it:} \emph{zavoyóvyvat}
  \item \emph{Say it:} \emph{zavoyóvyvat}, once more
  \item \emph{Say it:} a sentence of your own with \emph{zavoyóvyvat}
  \item \emph{Recall:} say the Russian for a revolution, then the Russian for an uprising, then say \emph{zavoyóvyvat} again
\end{itemize}

\begin{tcolorbox}[breakable,colback=teal!4,colframe=teal!35!black,title={Before you move on}]
What does \ru{завоёвывать} mean? (\textquotedblleft{}To conquer\textquotedblright{}.) Say it once more.
\end{tcolorbox}
\section[\texorpdfstring{\ru{править} (právit)}{править (právit)}]{\ru{править} (právit) — to rule, to govern; to correct}

\label{lesson:RU-C325-pravit}



\begin{quote}
Before the new one: say the Russian for an uprising, then the Russian for to conquer.
\end{quote}

\subsection*{You'll want to know: \ru{править}}
\textbf{\ru{править}} (\emph{právit}) — \textquotedblleft{}to rule, to govern; to correct\textquotedblright{}.

\begin{quote}
\ru{Царь} \ru{правил} \ru{страной} \ru{тридцать} \ru{лет}.
\end{quote}

\begin{quote}
\emph{Tsar právil stranóy trídtsat let.}
\end{quote}

\begin{quote}
The tsar ruled the country for thirty years.
\end{quote}

Imperfective. It takes the instrumental: \emph{\ru{править} \ru{страной}}. Also to correct a text: \emph{\ru{править} \ru{статью}}.

\begin{cousinweb}
\emph{\ru{править}} (EtymDB: \textquotedblleft{}to govern\textquotedblright{}) is built on \emph{\ru{правый}} \textquotedblleft{}right, straight\textquotedblright{}, an inherited Slavic word: to keep things straight.
\end{cousinweb}

\begin{grammarlens}[title={What you've built}]
One more word for telling, explaining and describing at B1.
\end{grammarlens}

\subsection*{Your turn}
\begin{itemize}
\raggedright
  \item \emph{Say it:} \emph{právit}
  \item \emph{Say it:} \emph{právit}, once more
  \item \emph{Say it:} a sentence of your own with \emph{právit}
  \item \emph{Recall:} say the Russian for an uprising, then the Russian for to conquer, then say \emph{právit} again
\end{itemize}

\begin{tcolorbox}[breakable,colback=teal!4,colframe=teal!35!black,title={Before you move on}]
What does \ru{править} mean? (\textquotedblleft{}To rule, to govern; to correct\textquotedblright{}.) Say it once more.
\end{tcolorbox}
